\chapter{关键算法优化}

第四章介绍的算法后端在结构、参数和对象规模上差异明显。本章围绕三个实际瓶颈展开：
Sageloli PARAM-I相关扰动的单个稠密矩阵理论存储约为1.18~TiB，SRGG参考对象存在较大的文本表示开销，Lai在固定身份对
通信中会反复生成相同叶密钥。三项优化分别改变计算组织、对象编码和运行时生命周期。

\section{Sageloli结构化环采样优化}

\subsection{稠密矩阵展开的资源瓶颈}

Sageloli的身份发钥和签名都需要在格陪集中采样。相关扰动由陷门决定其形状，再与离散采样
结合，生成满足方程和范数要求的短向量。环元素本来以多项式系数表示，乘法可以直接在
负循环环上计算；若把每个环元素展开为旋转矩阵，陷门和协方差会随之变成规模很大的实数矩阵。

PARAM-I取$d=2048$、$\ell=132$、$k=65$，签名相关扰动的系数维数为
\[
 D=d(\ell+k)=403456.
\]
仅显式保存一个$D\times D$双精度实数矩阵就需要约1.18~TiB，还没有计入矩阵分解和临时工作区。
实现采用结构化表示，直接计算所需的矩阵作用。

实现始终保留多项式结构。公钥、主陷门和身份陷门分别以$\ell\times d$、
$(\ell-k)\times k\times d$和$\ell\times k\times d$的系数数组保存，环乘法在原环上执行。
每次发钥和签名结束后仍以精确整数算术检查
\[
 A=[\bar A\mid g-\bar AR],\qquad
 AT_s=H_1(P,\rho,s)-g\pmod q,
\]


\subsection{逐频率低秩分解}

设当前陷门为$T$，其形状为$\ell\times k$个环元素，并令
\[
 J=\begin{bmatrix}-T\\I_k\end{bmatrix},\qquad
 a=s^2-r^2,\qquad c=9r^2.
\]
相关扰动所需协方差为
\[
 \Sigma=aI-cJJ^\star,
\]
其中$\star$表示环上的伴随。$aI$是标量矩阵，修正项由只有$k$列的$J$生成。对任何与
$\operatorname{span}(J)$正交的向量，$J^\star x=0$，所以修正项$JJ^\star$不起作用，平方根算子
在逐频率表示中作用于$J_\nu$的至多$k$维列空间。

首先对多项式执行负循环Fourier变换。在每个频率$\nu$上，环乘法变为复数乘法，$J$对应
$(\ell+k)\times k$复矩阵$J_\nu$。随后计算
\[
 G_\nu=J_\nu^*J_\nu
      =U_\nu\operatorname{diag}(\lambda_i)U_\nu^*.
\]
在PARAM-I中，每个频率分解一个$65\times65$的Gram矩阵。设$G_\nu$的单位特征向量为$v_i$、特征值为$\lambda_i$，则$J_\nu v_i/\sqrt{\lambda_i}$给出$\operatorname{span}(J_\nu)$中的正交方向，其对应参数值为$a-c\lambda_i$；在该子空间的正交补上，平方根作用为$\sqrt a$倍缩放。

\begin{figure}[H]
\centering
\begin{tikzpicture}[node distance=8mm and 7mm]
 \node[flow,text width=39mm] (poly) {陷门多项式$T$\\保存系数数组};
 \node[flow,text width=39mm,right=of poly] (fourier) {负循环Fourier变换\\$J_\nu$：$197\times65$};
 \node[flow,text width=39mm,right=of fourier] (gram) {计算$J_\nu^*J_\nu$\\分解$65\times65$矩阵};
 \node[check,text width=39mm,below=of gram] (apply) {对新随机向量应用\\低秩平方根};
 \node[check,text width=39mm,left=of apply] (inverse) {逆变换回系数域\\逐系数离散采样};
 \node[check,text width=39mm,left=of inverse] (verify) {检查精确模方程\\和原方案范数条件};
 \draw[flowarrow](poly)--(fourier);
 \draw[flowarrow](fourier)--(gram);
 \draw[flowarrow](gram)--(apply);
 \draw[flowarrow](apply)--(inverse);
 \draw[flowarrow](inverse)--(verify);
\end{tikzpicture}
\caption{Sageloli结构化相关扰动的计算流程（矩阵形状对应PARAM-I）}
\label{fig:ring-sampler}
\end{figure}

根据小矩阵的特征分解构造平方根算子
\[
 L_\nu=\sqrt a I-
 J_\nu U_\nu\operatorname{diag}\!\left(
 \frac{c}{\sqrt a+\sqrt{a-c\lambda_i}}
 \right)U_\nu^*J_\nu^*.
\]
该算子满足$L_\nu L_\nu^*=aI-cJ_\nu J_\nu^*$。应用算子时，程序先把新生成的标准随机向量
投影到$J_\nu$的列空间，再用$65\times65$特征分解修正这些分量；其余分量直接乘以$\sqrt a$。
实现以小矩阵与向量乘法直接计算$L_\nu$对随机向量的作用。逆变换回系数域后，继续执行原方案要求的
离散Gaussian和三进制陪集采样。该方法沿用格陷门采样中的扰动思想\cite{micciancio2012trapdoors}，
并针对当前方案的环结构组织计算。

主要存储量由陷门系数和逐频率小矩阵计划组成，规模为$O(d\ell k+dk^2)$。
固定陷门的系数数组、Fourier表示和各频率特征分解可以复用；每次发钥或签名仍重新生成Gaussian
源向量、相关扰动和陪集样本。表~\ref{tab:sageloli-representation}比较两种计算组织。

\begin{table}[H]
\centering\small
\caption{Sageloli稠密表示与结构化实现对照}
\label{tab:sageloli-representation}
\begin{tabularx}{\textwidth}{@{}>{\raggedright\arraybackslash}p{.2\textwidth}>{\raggedright\arraybackslash}p{.24\textwidth}>{\raggedright\arraybackslash}p{.25\textwidth}>{\raggedright\arraybackslash}X@{}}
\toprule
\rowcolor{black!8} 计算组织 & 核心分解规模 & 主要存储 & PARAM-I结果 \\
\midrule
系数空间稠密展开 & 一个$403456\times403456$实矩阵 & 单个双精度矩阵约1.18~TiB，尚未计分解工作区 & 理论存储估算 \\
逐频率低秩实现 & 2048个$197\times65$的$J_\nu$，各分解$65\times65$的Gram矩阵 & $O(d\ell k+dk^2)$系数与数值计划 & 完成发钥、签名和公开验证，实测峰值RSS约1.36~GiB \\
\bottomrule
\end{tabularx}
\end{table}

Sageloli原论文给出SIS／RSIS两类IBS构造、ROM／QROM分析以及Table~6参数\cite{sageloli2023ibs}。
本作品围绕其中RSIS参数补充系数原生对象、逐频率低秩SampleF、目标尺寸编解码和精确环方程检查，
并把后端接入通用IB-ME流程。

\subsection{数值采样与精确代数校验}

Fourier变换和特征分解使用数值运算，方案中的模方程则必须精确成立。实现将两类职责分开：
数值计划在建立时检查正定充分条件和变换误差；回到系数域后，以整数算术检查环方程。大模数
$q=3^k$的公开乘积使用FLINT整数多项式，私有小整数乘法采用双素数NTT与CRT恢复，并检查
整数卷积幅度处于可唯一恢复范围。

签名生成后继续执行原方案的范数界和验签方程。浮点计算负责生成相关扰动，精确算术负责确认
最终代数关系。第七章给出PARAM-I、PARAM-II的实测时间、峰值内存及验签结果。

\section{SRGG对象的范围感知编码}
\label{sec:srgg-codec}

\subsection{参考表示的空间开销}

SRGG的公钥、陷门和签名包含大量多项式系数。JSON会把每个整数写成十进制字符串，并重复保存
负号、逗号、括号和字段名。对于大对象，这些语法字符和文本数字会形成显著的额外开销。

紧凑编码依据各字段的数学取值范围选择位数。例如，取值为$-1,0,1$的陷门系数只有三种可能，
每项用2~bit即可表示；公开模$Q=15744063$的系数位于$[0,Q)$，每项使用24~bit。短密钥和
签名响应先以Zigzag映射处理符号，再选择能够容纳该字段的最小统一位宽。

\begin{table}[H]
\centering\small
\caption{SRGG不同字段的紧凑编码方法}
\label{tab:codec-fields}
\begin{tabularx}{\textwidth}{@{}p{.23\textwidth}p{.25\textwidth}X@{}}
\toprule
\rowcolor{black!8} 数据 & 存储方法 & 解码时检查 \\
\midrule
公开模$Q$系数 & 每系数24~bit & $0\leq x<Q$ \\
三值陷门系数 & 每系数2~bit & 拒绝第四个未使用码字 \\
短密钥、响应整数 & Zigzag与最小统一位宽 & 数值范围、位宽最小性 \\
二元挑战 & 每系数1~bit & 固定多项式次数与长度 \\
类型、参数与上下文 & 固定头部和长度前缀 & 对象类别、公钥和参数一致 \\
\bottomrule
\end{tabularx}
\end{table}

对包含$N$个系数、统一位宽为$b$的字段，主体长度为$\lceil Nb/8\rceil$字节，再加必要头部。
两种格式无损互转，签名与验证共用同一算法。

\subsection{四类对象的规范编解码}

范围感知编码同时用于公钥、机构状态、用户密钥和签名。四类对象使用独立类型标识，系数顺序固定，
末字节中未使用的位必须为零。解析器先根据参数和形状计算预期长度，再读取数据，并检查对象类型、
身份、公钥上下文和系数范围。解码后重新编码必须得到相同字节串，从而排除同一数学对象的多种
字节表示。

对象变短可以减少密钥文件和传输字节，但收益取决于对象在完整协议中的占比。例如，IBS签名
变短会直接减少对称加密载荷，而HIBE种子密文和一次性认证仍构成固定开销。第七章同时报告四类对象
的JSON与紧凑格式长度、编解码时间以及完整密文组成，以区分对象级收益和协议级收益。

SRGG原论文给出PARAM-I、II、III三档规范、QROM安全分析和概念验证实现\cite{jia2025ibs}。
本作品在此基础上为公钥、机构状态、用户密钥和签名建立统一的无损规范编码，补充严格解析、状态
持久化、通用IB-ME接入及完整载荷核算。表~\ref{tab:srgg-protocol-impact}按两份PARAM-I功能运行记录核算签名与完整密文长度。两次运行分别建立系统，采用相同算法参数、身份长度和消息长度，签名分别采用参考与紧凑格式。

\begin{table}[H]
\centering\small
\caption{两次LATTE＋SRGG PARAM-I运行的签名与密文字节核算}
\label{tab:srgg-protocol-impact}
\begin{tabular}{@{}lrrr@{}}
\toprule
\rowcolor{black!8} 对象 & 参考表示 & 紧凑表示 & 减少 \\
\midrule
% BEGIN DATA srgg-protocol-impact.tex
SRGG签名 & 35,289~B & 13,400~B & 21,889~B（62.0\%） \\
完整B密文 & 171,114~B & 149,231~B & 21,883~B（12.8\%） \\
% END DATA srgg-protocol-impact.tex
\bottomrule
\end{tabular}
\begin{minipage}{.96\textwidth}\footnotesize
来源：LATTE＋SRGG PARAM-I两次独立功能运行，消息10~B；公钥、主私钥与用户密钥采用JSON格式，签名采用表中对应格式。

\end{minipage}
\end{table}

\section{Lai接收会话优化}
\label{sec:lai-session}

\subsection{固定身份对下的重复委派}

在Lai组合中，接收者持有第一层密钥$\mathsf{dk}_r$。解密来自发送者$s$的消息前，需要运行
\[
 \mathsf{dk}_{r,s}\leftarrow\mathsf{Derive}(\mathsf{dk}_r,s).
\]
同一接收者连续读取同一发送者的消息时，委派输入保持不变，逐消息生成叶密钥会重复支付相同的
计算成本。与此同时，每条消息的HIBE种子密文、对称密文和IBS签名都不同，仍须逐条认证与解密。

作品引入显式接收会话：会话建立时生成并验证一次叶密钥，随后在固定双方身份下复用。图
\ref{fig:receiver-session}对比两种计算安排。

\begin{figure}[H]
\centering
\begin{tikzpicture}[node distance=6mm and 6mm]
 \node[font=\small,anchor=east] (old) {逐条委派};
 \node[flow,right=of old,text width=30mm] (o1) {委派＋解密1};
 \node[flow,right=of o1,text width=30mm] (o2) {委派＋解密2};
 \node[flow,right=of o2,text width=30mm] (o3) {委派＋解密3};
 \node[font=\small,anchor=east,below=22mm of old] (new) {接收会话};
 \node[check,right=of new,text width=30mm] (prep) {委派一次\\验证并保留叶密钥};
 \node[check,right=of prep,text width=30mm] (dec) {解密1、2、3\\逐条认证与验签};
 \node[check,right=of dec,text width=30mm] (close) {关闭会话\\停止使用叶密钥};
 \draw[flowarrow](o1)--(o2);
 \draw[flowarrow](o2)--(o3);
 \draw[flowarrow](prep)--(dec);
 \draw[flowarrow](dec)--(close);
\end{tikzpicture}
\caption{Lai接收会话将与密文无关的委派移至准备阶段}
\label{fig:receiver-session}
\end{figure}

\subsection{身份绑定与会话生命周期}

会话建立时先检查公开上下文、接收者身份和一级密钥绑定；生成叶密钥后再检查完整$(r,s)$路径及
陷门方程。会话只服务于当前公钥和身份对，关闭后停止解密。每条消息继续执行规范解析、外层字节
认证、HIBE解封装、AES-GCM验证、IBS验签和一次性公钥承诺核对。

会话叶密钥保留在进程内，并绑定当前公钥及双方身份；长期接收方密钥采用持久化格式保存。每条消息独立生成随机量并执行完整认证。

\subsection{连续消息的摊销成本}

\textbf{会话内单条解密由526秒缩短到28.6秒，提速约18.4倍，耗时减少约94.6\%。}
每条消息减少约498秒的接收等待，主要来自复用已经生成的叶密钥。接收会话的一次准备时间约为510秒。

设准备时间为$T_{\rm prep}$，会话内单条解密时间为$T_{\rm sess}$，逐条委派路径的单条时间为
$T_{\rm original}$。处理$N$条消息时，总成本和平均成本分别为
\[
 T_{\rm total}(N)=T_{\rm prep}+NT_{\rm sess},\qquad
 \bar T(N)=T_{\rm prep}/N+T_{\rm sess}.
\]
代入实测阶段中位数$T_{\rm prep}=509.548$秒、$T_{\rm original}=526.272$秒、
$T_{\rm sess}=28.568$秒，可得表~\ref{tab:session-amortization}。

\begin{table}[H]
\centering\small
\caption{按实测阶段时间计算的累计解密成本}
\label{tab:session-amortization}
\begin{tabular}{@{}rrrr@{}}
\toprule
\rowcolor{black!8} 消息数 & 逐条委派／分钟 & 接收会话／分钟 & 累计耗时减少 \\
\midrule
2 & 17.5 & 9.44 & 46.2\% \\
10 & 87.7 & 13.3 & 84.9\% \\
100 & 877 & 56.1 & 93.6\% \\
\bottomrule
\end{tabular}
\par\smallskip\begin{minipage}{.96\textwidth}\footnotesize
成本模型：逐条委派为$NT_{\rm original}$；接收会话包含一次准备，计为$T_{\rm prep}+NT_{\rm sess}$。
\end{minipage}
\end{table}

累计成本低于原路径的条件为
\[
 N>\frac{T_{\rm prep}}{T_{\rm original}-T_{\rm sess}}=1.024.
\]
因此，接收第二条消息时即可摊回准备开销。处理100条周期报送消息的模型总时间由约14.6小时
缩短至约56分钟。第七章列出各阶段的测量次数与会话叶密钥的存储长度。

\section{本章小结}

三项优化分别处理大型相关扰动计算、对象编码和固定身份对下的重复委派。Sageloli采用逐频率低秩分解；SRGG按字段范围进行规范二进制编码；Lai通过绑定身份与公钥的会话复用叶密钥。第七章给出相应的时间、空间与通信量结果。
